\documentclass[12pt, a4paper]{report}

\usepackage[utf8]{inputenc}
\usepackage{geometry}
\geometry{a4paper, margin=1in}
\usepackage{hyperref}
\usepackage{graphicx}
\usepackage{listings}
\usepackage{xcolor}
\usepackage{titlesec}
\usepackage{booktabs}
\usepackage{amsmath}

% Formatting for code snippets
\definecolor{codegray}{rgb}{0.5,0.5,0.5}
\definecolor{codepurple}{rgb}{0.58,0,0.82}
\definecolor{backcolour}{rgb}{0.96,0.96,0.94}
\definecolor{codeblue}{rgb}{0.13,0.13,0.8}

\lstdefinestyle{mystyle}{
    backgroundcolor=\color{backcolour},   
    commentstyle=\color{codegray}\itshape,
    keywordstyle=\color{codeblue}\bfseries,
    numberstyle=\tiny\color{codegray},
    stringstyle=\color{codepurple},
    basicstyle=\ttfamily\footnotesize,
    breakatwhitespace=false,         
    breaklines=true,                 
    captionpos=b,                    
    keepspaces=true,                 
    numbers=left,                    
    numbersep=5pt,                  
    showspaces=false,                
    showstringspaces=false,
    showtabs=false,                  
    tabsize=2
}
\lstset{style=mystyle}

\title{\Huge{\textbf{Comprehensive Project Development Report}} \\ \vspace{0.5cm} \Large{\textbf{Full-Stack Room Booking System}}}
\author{Mokshith Kumar Reddy \& AI Assistant}
\date{\today}

\begin{document}

\maketitle

\begin{abstract}
This document serves as a comprehensive, step-by-step technical report detailing the complete lifecycle of the Room Booking System development. It covers the initial requirements gathering, technology stack selection, database design, full-stack implementation, advanced algorithmic logic for room availability, and the final deployment to GitHub. Every major step taken by the AI assistant in collaboration with the user is documented here in extreme detail to provide a full understanding of the project's internal workings.
\end{abstract}

\tableofcontents
\newpage

\chapter{Project Inception \& Requirements}

\section{The Objective}
The primary objective was to build a full-stack web application tailored for an academic/corporate environment (specifically handling blocks like LHC, CSE, A Block, etc. at IIT Hyderabad). The system needed to allow authorized members to book rooms without conflicts, while an admin oversees all operations.

\section{Core Requirements \& Rules}
Based on the initial prompts provided, the strict rules for the system were:
\begin{itemize}
    \item \textbf{No Public Sign-Up:} Members cannot create their own accounts.
    \item \textbf{Admin Control:} A designated Admin creates all member accounts and provides them with login credentials.
    \item \textbf{Language \& Tech:} Strict use of JavaScript, React, Tailwind CSS, Express.js, PostgreSQL, and Prisma.
    \item \textbf{Overlap Logic:} The system must rigorously check date and time overlaps to prevent double booking.
    \item \textbf{UI Aesthetics:} The design must be premium, utilizing dark mode and glassmorphism.
\end{itemize}

\chapter{Database Design \& Architecture}

\section{Choosing PostgreSQL \& Prisma}
We opted for PostgreSQL for its robustness in handling relational data (Users $\rightarrow$ Bookings $\leftarrow$ Rooms). Prisma was selected as the Object-Relational Mapper (ORM) because of its strict type safety and auto-generated migrations.

\section{The Schema Design}
The database was structured into four primary models in \texttt{schema.prisma}:
\begin{enumerate}
    \item \textbf{User Model:} Fields include \texttt{id, name, email, password}, and an Enum for \texttt{role} (\texttt{ADMIN} or \texttt{MEMBER}).
    \item \textbf{Block Model:} Fields include \texttt{id, blockName}. This allows rooms to be grouped logically.
    \item \textbf{Room Model:} Fields include \texttt{id, roomName, capacity, availabilityStatus}, and a foreign key linking to a \texttt{Block}.
    \item \textbf{Booking Model:} The most complex model linking a \texttt{User} to a \texttt{Room}. It tracks the \texttt{date}, \texttt{startTime}, \texttt{endTime}, \texttt{purpose} (OA, Interview, PPT), and \texttt{status}.
\end{enumerate}

\chapter{Backend API Implementation}

\section{Authentication Flow (JWT \& bcrypt)}
Security is paramount since anyone accessing the system must be pre-authorized. We utilized \texttt{bcrypt} to hash passwords before saving them to the database. Upon a successful login via the \texttt{/api/auth/login} route, the server generates a JSON Web Token (JWT) signed with a secure secret. 

Two custom middleware functions were created:
\begin{itemize}
    \item \texttt{authenticateToken}: Extracts the JWT from the HTTP authorization header, verifies it, and attaches the user data to the request.
    \item \texttt{authorizeRole('ADMIN')}: Checks if the attached user has the \texttt{ADMIN} role. If a regular member attempts to hit an admin route, they receive a 403 Forbidden response.
\end{itemize}

\section{The Core Overlap Algorithm}
The most technically challenging prompt was implementing the room search logic. The user requested that the search parameters include date, start time, end time, purpose, required capacity, and an optional block.

To determine if a room is available, we must check if any existing booking for that room on that date overlaps with the requested time. The mathematical logic implemented is:

\begin{lstlisting}[language=Java, caption=Overlap Logic Implementation]
const hasOverlap = room.bookings.some((booking) => {
    // There is NO overlap if the existing booking ends before/exactly when the new one starts
    // OR if the existing booking starts after/exactly when the new one ends.
    const noOverlap = (booking.endTime <= requestedStartTime) || (booking.startTime >= requestedEndTime);
    
    // If noOverlap is false, it means they DO overlap.
    return !noOverlap; 
});
\end{lstlisting}

\section{Refining the Search UX}
In a subsequent prompt, the user requested an enhancement: instead of filtering out unavailable rooms completely from the search results, the system should show them but tag them as unavailable. 

To achieve this, the backend \texttt{searchRooms} controller was modified to map over \emph{all} rooms that meet the capacity and block criteria, evaluating the overlap logic, and attaching an \texttt{isAvailableForBooking} boolean flag to the response payload.

\chapter{Frontend Development}

\section{Vite \& React Configuration}
The client-side was scaffolded using Vite for optimal performance. React Router DOM was implemented to manage client-side routing, dividing the app into public routes (Login) and protected routes (Dashboard, Search Rooms, Manage Users).

\section{Premium UI Design}
As requested, the UI was designed to be visually stunning. Tailwind CSS was heavily utilized to create a dark-themed environment. We implemented ``Glassmorphism'' by using translucent backgrounds with background-blur utilities (e.g., \texttt{bg-white/5 backdrop-blur-lg border border-white/10}). This gave the application a highly modern, sleek appearance.

\section{Frontend Integration}
Axios was configured with an interceptor to automatically attach the JWT token (stored in \texttt{localStorage}) to every outgoing request. This ensured seamless communication with the protected Express APIs.

When rendering the room search results, the frontend consumes the \texttt{isAvailableForBooking} flag to conditionally render a green "Available" badge or a red "Unavailable" badge. Furthermore, if a room is unavailable, the ``Book Room'' button is disabled and its opacity is lowered.

\chapter{Data Extraction \& Seeding}

\section{Parsing the PDF Document}
The user provided a PDF document (\texttt{OCS\_IITH\_Full\_Task\_Document.pdf}) containing a reference table of exact rooms and capacities (e.g., \texttt{A-Class Room 320}, \texttt{LHC-01}). To make the application realistic, we utilized a background terminal command to extract the text from this PDF using \texttt{pdftotext}.

\section{The Seed Script}
A \texttt{seed.js} script was engineered to automatically populate the PostgreSQL database. The script:
\begin{enumerate}
    \item Hashes passwords for demo accounts.
    \item Creates an \texttt{ADMIN} user (\texttt{admin@roombooking.com}).
    \item Creates a temporary \texttt{MEMBER} user (\texttt{tempuser@roombooking.com}).
    \item Iterates through an array of 10 blocks and inserts them.
    \item Iterates through 33 specific rooms extracted from the PDF, associating them with the correct block IDs and capacities.
\end{enumerate}

\chapter{Deployment \& Environment Management}

\section{Port Conflict Resolution}
During the local testing phase, we encountered an \texttt{EADDRINUSE} error indicating that port \texttt{5000} was occupied (a common issue on macOS due to AirPlay Receiver). We resolved this by modifying the \texttt{.env} file to use port \texttt{5001}, updating the Vite proxy configuration to target \texttt{http://localhost:5001}, and forcefully killing the orphaned background processes.

\section{GitHub CLI Integration}
The final request was to push the entire project to a new GitHub repository named \texttt{OCS-Projects}. Because the GitHub CLI (\texttt{gh}) requires interactive browser authentication, an AppleScript command was executed to pop open a native macOS Terminal window directly on the user's screen. This allowed the user to securely authenticate via OAuth. 

Once authenticated, the following commands were executed to successfully push the codebase:
\begin{lstlisting}[language=bash]
git init
git add .
git commit -m "Initial commit: Full-stack Room Booking System"
gh repo create OCS-Projects --public --source=. --remote=origin --push
\end{lstlisting}

\chapter{Conclusion}
This project successfully transformed a raw set of requirements into a production-ready, full-stack room booking architecture. By systematically breaking down the database modeling, API creation, advanced mathematical scheduling logic, UI/UX design, and deployment processes, we constructed a highly secure and visually impressive platform. The active collaboration between the user's iterative prompts and the AI's technical execution resulted in a flawless final product.

\end{document}
